% THIS IS A LATEX TEMPLATE FILE FOR PAPERS INCLUDED IN THE
% *Anthology of Computers and the Humanities*.
\documentclass[final]{anthology-ch}

% LOAD LaTeX PACKAGES
\usepackage{booktabs}
\usepackage{graphicx}

% TITLE OF THE SUBMISSION
\title{Understanding Ephemerality and Loss in Digital Humanities Projects: What Surveys of the Field Can Reveal}

% AUTHOR AND AFFILIATION INFORMATION
\author[1]{Kathryn C. Wymer}[
  orcid=0000-0003-0100-2692
]

\affiliation{1}{University of North Carolina at Chapel Hill, Chapel Hill, NC, USA}

% KEYWORDS
\keywords[eng]{digital humanities, sustainability, project ephemerality, graceful degradation, survey research, digital preservation}

% METADATA FOR THE PUBLICATION
\pubyear{2025}
\pubvolume{5}
\pagestart{70}
\pageend{76}
\paperorder{9}
\conferencename{From Manuscript to Megabyte: Building Sustainable and Open Futures in the Digital Humanities}
\conferenceeditors{James B. Harr III, Emma Stanley, Jenna Rinalducci, and Donna Kain}
\doi{10.63744/0zIOrMmzCwMg}
\customhead{}

\addbibresource{bibliography.bib}

%%%%%%%%%%%%%%%%%%%%%%%%%%%%%%%%%%%%%%%%%%%%%%%%%%%%%%%%%%%%%%%%%%%%%%%%%%%
% HERE IS THE START OF THE TEXT
\begin{document}

\maketitle

\nocite{arneil2023planning,otis2023follow}

\begin{abstract}
This paper discusses the ever-changing technology power the digital humanities and the resulting ephemerality and loss of DH projects, specifically looking at surveys of the field for insight. The Graceful Degradation Survey, originally implemented in 2009 and conducted again in 2022, intended to learn more about how DH projects decline or end and how practitioners plan for or react to such outcomes. Topics addressed in this discussion include the longitudinal nature of this data, which is also worth comparing to the Endings Project at the University of Victoria and to developments in DH research. Some findings include data about the demographics of the field (which have remained somewhat static over time); emergent understandings about decline but slow implementation of solutions; concerns about what might be the best practices for maintaining or ending projects; and concerns about whether planning for project ends has the potential to stifle innovation or creativity.
\end{abstract}

\section{Introduction: DH as Research on Fire} \label{sec:intro}

In \enquote{Platforms are Infrastructures on Fire,} Paul N. Edwards \cite[p.~313]{edwards2021platforms} wrote that networked software platforms \enquote{are fast, but they're flammable,} noting that new platforms can emerge quickly but can also vanish just as quickly. As an example, Edwards reminds us that the now-defunct social media platform Friendster had 115 million users at one point in the not-too-distant past. It's easy for anyone who has used social media for more than a few years to think of other services that have either disappeared or nearly faded from use: for example, Vine, MySpace, or GooglePlus. New ways of connecting online appear all the time, and once-popular platforms fade away. Other platforms change in ways that create significant impacts, such as Twitter's metamorphosis into X. These kinds of transitions may be cumbersome, but they are commonplace. The rapid pace of this type of technology to rise and fall is to be expected. It's the way the systems work.

For similar reasons, it seems logical to argue that the digital humanities in many ways can be thought of as research on fire. Traditional publications in books and journals can take years to research and write and then spend additional months to years in the review process. Once published, these works are meant to last indefinitely. Whether these books or journals are used indefinitely (or whether they are just gathering dust on a shelf) is not a metric many scholars can produce.\footnote{Naturally, citations provide metrics, but we have little way to know when a work is used in a classroom or simply read by an interested scholar.} Nevertheless, there is a sense of permanence to traditional forms of scholarship that does not exist for the digital humanities.

How many DH projects can you think of that no longer exist? Or that don't exist in the same format as they used to? DH projects take shape and often burn out as quickly as do the software platforms with which they co-exist. In fact, sometimes changes to platforms disrupt the DH projects that they are built on (with the deprecation of Adobe Flash or functional changes to Twitter/X being memorable examples with widespread impacts).

\section{Do We Need DH Projects to Last?} \label{sec:need}

If impermanence in the digital world should really be considered an intentional feature rather than a flaw, why are scholars worried about the impacts of impermanence in DH? First and foremost for many university scholars, the academy has had difficulty in learning how to value digital work. Traditional forms of scholarship with their connotation of permanence have long been the standard for how to evaluate a scholar's relevance. A digital project that begins and ends quickly poses a challenge for many who attempt to evaluate scholarship. Although organizations like the Association for Computers and the Humanities \cite{ach2016guidelines} have created guidelines on how to value DH work without respect to its longevity, institutional responses to ephemeral digital projects continue to vary greatly. Some DH practitioners therefore have assumed that their projects must last indefinitely to be relevant to other scholars. That attitude leads to a lack of planning for a project ending, leads to often incorrect assumptions that continued maintenance will be available, and sometimes results in eventual \enquote{Zombie DH} projects that are neither concluded nor ongoing and that have broken links or other nonfunctioning elements.\footnote{This term comes from the work of Quinn Dombrowski \cite{dombrowski2018zombie}, who used it to describe the process of moving from maintaining projects, to letting them lapse, and then perhaps to re-discovering them.}

A corollary problem is that DH scholars often fail to distinguish the parts of their project that have lasting importance from those parts that could be brought to a close \cite[p.~11]{wymer2021introduction}. For instance, a project's data might be worth keeping and have lasting usefulness even outside of an interactive digital project space that may not be maintainable. Without plans for how to extract, store, or share it, that useful data may be lost along with the rest of the project. Alternately, without a clear project management plan or data management plan, a DH project could linger on an unmaintained internet site where it might provide uncurated information that could prove misleading or confusing to users.

\section{What Surveys of the DH Field Can Tell Us about Ephemerality and Loss} \label{sec:surveys}

\subsection{Surveys under Discussion}

Understanding how scholars have come to deal with ephemerality and loss in digital humanities projects was the motivation behind the Graceful Degradation survey project, which was originally administered in 2009. Since that time, much work on the topic has been done, often in the context of sustainability studies. Because of the continued interest in how and why DH projects end, the Graceful Degradation project was followed up with a second survey in 2022 \cite{wymer2026graceful}. An advantage of this survey process is the longitudinal data it gathered about the practical landscape of DH projects in the intervening years. Whereas great strides have been made in theorizing DH and considering the impact of technological advances on the arts and humanities, it is worth investigating how much DH practitioners have been able to adapt their scholarship to rapidly evolving technological changes. Comparing the results of the two surveys can help us understand where DH has progressed but also what steps DH practitioners may need to take in the coming years. Notably, this work overlaps with another major effort to understand how and why DH projects should conclude: The Endings Project at the University of Victoria \cite{endings2024about, holmes2023tamagotchis}. Their 2017--2018 survey provides an informational midpoint for comparison, and it is useful to assess their recommendations for how to conclude or preserve DH projects \cite{goddard2023whats, endings2024principles} with respect to the responses to the Graceful Degradation surveys.

\subsection{Notable Differences in the Surveys Affecting Comparison of Data}

\begin{itemize}
  \item Timeframes: 2009, 2017--2018, and 2022
  \item Different originators with different location base (US/UK, Canada, US)
  \item Different questions
  \item Different methods of distributing survey (though all online)
\end{itemize}

The origins of the Graceful Degradation surveys date back to work done in 2009--2010 by Bethany Nowviskie and Dot Porter. They created the first Graceful Degradation survey in order to see how many digital projects decline and the nature of that decline. One of the goals was to make recommendations about improving the processes for how to deal with project decline. Nowviskie and Porter \cite{nowviskie2010graceful} described the survey methods for a presentation at the Digital Humanities 2010 conference at King's College London. In that presentation, which is publicly available, they reported initial findings from the survey. They also privately shared a fuller data set with Kathryn Wymer in order to develop a subsequent survey and compare the two.

The 2022 Graceful Degradation Survey was conceived as an extension of the original survey. Using the same basic questions, the idea was to compare how the situation of DH project decline has developed. To the casual observer, they would appear as essentially the same survey. However, in the interest of transparency, it is true the two surveys are not exactly the same.

\subsection{What These Surveys Have in Common}

\begin{itemize}
  \item An intended purpose of understanding how DH practitioners deal with digital projects that end.
  \item An anticipated goal of collecting data to help back solutions or justifications for how to deal with sustainability or preservation of projects.
  \item A very similar number of respondents and comparable (though not identical) questions.
\end{itemize}

The Endings Project had 127 respondents \cite{endings2024survey}. For the Graceful Degradation Surveys, there were similar response rates, with 102 responses in 2009 and 118 responses in 2022. Participants were both not required to complete all questions, so total responses to each question varied. Most respondents answered the general questions about DH and project decline. Fewer responded to additional detailed questions about a project they had worked on, with some of those questions receiving about a 50\% response rate.

\subsection{Several Key Takeaways}

\subsubsection{Demographics of the DH Field Remain Stable}

One notable consistency between the two graceful degradation surveys is that the majority of respondents said that they had most recently completed DH work while affiliated with a university with a research emphasis. In 2009, 78\%, and in 2022, 77\% of those answering the question identified as such. Those affiliated with a college or university with a teaching emphasis were 14\% and 13\% respectively. At these institutions, respondents had a wide variety of faculty, staff, and administrative roles, but only 12 respondents in 2009 and 15 in 2022 identified as graduate students or postdocs. Considering how much time elapsed between surveys and the different methods of disseminating information in those periods, this consistency is remarkable and suggests something static about participation in DH.

\subsubsection{The 2022 Survey Captured New Voices, but Only Certain Perspectives}

\paragraph{Who was not represented in the survey?}
One group not represented in any of these surveys is people who work in pre-university education (designated as K--12 in the United States). Though more research traditionally is done at the university level, there are many pedagogical opportunities for DH in the K--12 setting. Though it is not clear, it seems reasonable to guess that some of the survey respondents from 2022 were representing university coursework and pedagogy as part of their DH projects. Any K--12 pedagogical work would fit here as well. It is possible that there is little DH work currently being done in K--12,\footnote{One project intended to promote DH work in K--12 is the work of Kelly Hammond \cite{hammond2022k12}.} but it is more likely that what work is being done is inadequately integrated into the larger conversation.\footnote{There are certainly some notable examples of DH literature on K--12 education, though they are less common than discussions of undergraduate and postgraduate work. For instance, work done by precollege students in partnership with a university was described by Cooke and Famiglietti \cite{cooke2023k12dh}.}

It is also the case that there were no responses from individuals who identified as undergraduate students in either 2009 or 2022. Though it is certain that some DH projects are happening in undergraduate classrooms or with the participation of undergraduate researchers, data from the perspectives of those students was not captured. Since it is also possible that survey respondents have an unclear view of \enquote{what counts} as DH, some may have declined to include small pedagogical projects, and those therefore may be underrepresented. In that case, the Principal Investigator for such projects may not have considered inviting undergraduates to respond to the Graceful Degradation Survey. Likewise, it seems probable that undergraduate voices are underrepresented in conversations about DH more generally.

\subsubsection{Implementations of Solutions to Problems Are Slow}

In the answers to the Graceful Degradation survey, there were lots of discussions about the importance of data management plans. In response to the question \enquote{Do you feel that the views or practices of your local institution or the larger academic community have evolved in response to transitional phases for projects, or in response to the decline of digital humanities projects?} there was a strong discussion of sustainability and maintenance. In the intervening years, the Endings Project had begun making recommendations on considering how to close down and archive projects, and that kind of move in DH is reflected in 2022 survey answers. However, some responders still noted their belief and hope that projects could or should be maintained indefinitely, often in libraries, as reflected in this answer:

\begin{quote}
My library---which one would think would be dedicated to the acquisition, preservation, and dissemination of knowledge---has refused to fulfill these responsibilities in regard to the websites we created to preserve and disseminate our work.
\end{quote}

For those respondents who thought more conversation was happening around the end of projects, there was frequent mention of the inadequacy of those conversations in their answers about whether practices have evolved:

\begin{itemize}
  \item \enquote{Keeping lots of orphaned projects online is expensive and time consuming. The financial aspect does get the attention of administrators.. more so than the decline of an individual project or projects. But there continues to be reluctance to fully tackle the problem at its origin, i.e.\ when new projects are being planned\ldots}
  \item \enquote{Perhaps partially. There is a greater awareness about some of the issues faced by digital projects, such as long-term sustainability, but this does not mean that the institutional policies improved or that a sensible plan is in place}
  \item \enquote{I don't think so. We are very \enquote{now} focused}
  \item \enquote{Not really. We are asked to both maintain projects for the long term and take on new work, and those are at odds with each other. Faculty generally don't have a good understanding that a project is not a monograph and doesn't live on a shelf; it takes a lot of effort, money, and institutional knowledge to keep older projects running.}
  \item \enquote{I think these questions are attracting greater attention but not sustained conversation, at least in my field. There's a lot of attention to how to begin project and secure funding, but not much about the end or how to avoid decline.}
  \item \enquote{I don't think anyone knows what to do really.}
\end{itemize}

Based on the frequency of these kinds of responses, it seems work like that done by the Endings Project \cite{endings2024principles} could be usefully distributed more widely.

\subsubsection{Many Express Concerns about Best Practices for Maintaining or Ending Projects}

It is worth noting, however, that a number of respondents felt that certain kinds of project planning made them feel as if innovative projects were less welcomed:

\begin{itemize}
  \item \enquote{Sadly I think they have evolved in the wrong direction: there seems to be a \enquote{one size fits all} approach in a lot of library and/or IT departments, e.g.\ meaning they will provide a delivery platform in which every \enquote{end-of-life} DH project is thrown, but it fails to support their particular research foci and is thus meaningless. Since every DH project wants to be unique (this is after all the reason why it deserves funding), there is a mismatch of what is supported and what is expected from the researchers' point of view, and any true innovation is discouraged.}
  \item \enquote{Yes, I think the institution has become more wary and less inclined to support imaginative projects, which are seen as \enquote{boutique} and problematic to support.}
\end{itemize}

Some respondents to the question of whether their individual practices had changed said they were now \enquote{wary,} \enquote{cynical,} and more \enquote{modest} in their DH goals. There is a tension here between efficiency and creativity.

\subsubsection{Ongoing Debate: Does Planning for the End Stifle Creativity?}

Constance Crompton \cite{crompton2023boutique} offers insight on how DH practitioners might achieve a \enquote{marriage of sustainability and innovation,} but it is clear that it needs to be addressed at the individual project level. Perhaps one way of addressing this is to create spaces in which it is acceptable to experiment knowing that what one respondent described as \enquote{special snowflake} projects do have value even as there are requirements in place that such projects should be conducted with an open acknowledgement that there can be no expectation of long-term maintenance. It is also worth considering the context of minimal computing studies in these conversations \cite{risam2022introduction}.

\section{Final Thoughts} \label{sec:final}

One response to the 2022 survey summed up a good attitude toward project decline that they had developed as a result of their DH work: \enquote{We think that digital things are meant to last forever, but they rarely do. Even more so, decline and obsolescence should be built into projects because it's natural and going to happen. As a result of this experience I have come to believe that it's ok for projects to decline and/or fail.} This kind of attitude embraces the ephemerality of digital platforms and is perhaps the ideal for the larger DH community of practitioners as we move forward.

\printbibliography

\end{document}
